% Start preamble
\documentclass[12pt,a4paper]{article}
\usepackage{geometry}
 \geometry{
 a4paper,
 total={170mm,257mm},
 left=20mm,
 top=20mm,
 }
\usepackage[utf8]{inputenc}
\usepackage[T1]{fontenc}
\usepackage[pdftex]{graphicx}
\graphicspath{{./}}
\usepackage{enumitem}
\usepackage{pdfpages}
\usepackage{hyperref}
\usepackage{tikz}
\usepackage{attachfile}
\usepackage{epstopdf}
\usepackage{array}
\usepackage{multirow}
\usepackage{multicol}
\usepackage{float}
%\usepackage[table]{xcolor,colorbl}
\setlength{\textwidth}{16cm}
\setlength{\oddsidemargin}{-0.5cm}
\setlength{\evensidemargin}{-0.5cm}
%\setlenght{\headsep}{0cm}
\setlength\parindent{0pt}
%\setlength{\extrarowheight}{3pt}
\usepackage{listings}
%\usepackage{xcolor}

\input{arduinoLanguage.tex}
%%%%%% Counting oppgaves %%%%%%
 \newcount\questnum \questnum=0
 \def\oppgave{
            \advance\questnum by 1
	    \ifthenelse{\questnum>0\AND \questnum<9}
	    {
                \vskip 1cm
		\textbf{Oppgave}\hskip 5pt\the\questnum \hfill \hfill(6p)
		\vskip 3pt
		\hrule
	\vskip 0.5cm}
	{
                \vskip 1cm
		\textbf{Oppgave}\hskip 5pt \the\questnum \hfill \hfill(12p)
		\vskip 3pt \hrule \vskip 0.5cm }

		}

% End preamble

\begin{document}
\title{TIF Elektroteknikk 01}
\author{Faglærer: Fred-Olav Mosdal 90507684\\}
\maketitle
\newpage
\
\newpage
\oppgave{}%1
\vskip 2.5pt 
Hvilken forkortelse brukes for vekselstrøm og likestrøm?
\vskip 5pt 
\begin{tikzpicture}
	\draw[step=0.5cm,gray!20,very thin]  grid (17,6) ;
\end{tikzpicture}
\vskip 2.5pt 
\newpage

\oppgave{}%2
\vskip 2.5pt 
Lag en tabell og sett inn symbol og målenhet for strøm, spenning og resistans. 
\vskip 5pt 
\begin{tikzpicture}
	\draw[step=0.5cm,gray!20,very thin]  grid (17,6) ;
\end{tikzpicture}
\vskip 2.5pt 
\newpage
\oppgave{}%3
\vskip 2.5pt 
Hva er det som får elektronene i en ledning til å bevege seg?
\vskip 5pt 
\begin{tikzpicture}
	\draw[step=0.5cm,gray!20,very thin]  grid (17,6) ;
\end{tikzpicture}
\vskip 2.5pt 
\newpage
\oppgave{}%4
\vskip 2.5pt 
Tegn en krets med batteri, lyspære og ampermeter. Ampermetere skal måle strømmen som går i lyspæren
\vskip 5pt 
\begin{tikzpicture}
	\draw[step=0.5cm,gray!20,very thin]  grid (17,6) ;
\end{tikzpicture}
\vskip 2.5pt 
\newpage
\oppgave{}%5
\vskip 2.5pt 
Hvor stor resistans er det i en panelovn med U=230V og I=5A?
\vskip 5pt 
\begin{tikzpicture}
	\draw[step=0.5cm,gray!20,very thin]  grid (17,6) ;
\end{tikzpicture}
\vskip 2.5pt 
\newpage
\oppgave{}%6
\vskip 2.5pt 
Hvor stor resistans er det i en 15 meter lang tolederkabel av kobber som har et tversnitt på 1.5mm²?
\vskip 5pt 
\begin{tikzpicture}
	\draw[step=0.5cm,gray!20,very thin]  grid (17,6) ;
\end{tikzpicture}
\vskip 2.5pt 
\newpage
\oppgave{}%7
\vskip 2.5pt 
En parallellkobling består av tre resistanser, $R_1=30\Omega, R_2=45\Omega og R_3=60\Omega$. Hvor stor er totalresistansen. 
\vskip 5pt 
\begin{tikzpicture}
	\draw[step=0.5cm,gray!20,very thin]  grid (17,6) ;
\end{tikzpicture}
\vskip 2.5pt 
\newpage
\oppgave{}%8
\vskip 2.5pt 
Hvor stor effekt har en 12 volts lyspære når det går 0.5 ampere igjennomg den?
\vskip 5pt 
\begin{tikzpicture}
	\draw[step=0.5cm,gray!20,very thin]  grid (17,6) ;
\end{tikzpicture}
\vskip 2.5pt 
\newpage
\oppgave{}%9
\vskip 2.5pt 
Tre resistorer på $5\Omega, 10\Omega og 15\Omega$ er koblet i serie med en spenningskilde på 230V
\begin{enumerate}
	\item Tegn et skjema over koblingen
	\item Regn ut den totale resistansen. 
	\item Regn ut strømmen i kretsen
	\item Regn ut spenningen over motstanden på $10\Omega$
	\item Regn ut effekten i den samme motstanden
\end{enumerate}

\vskip 5pt 
\begin{tikzpicture}
	\draw[step=0.5cm,gray!20,very thin]  grid (17,6) ;
\end{tikzpicture}
\vskip 2.5pt 
\newpage
\end{document}
